% Section 4: the user's recovery outline, using the running batch-failure example.
\begin{tikzpicture}
\path[use as bounding box] (-.2,-1.35) rectangle (24.3,17.1);
\node[heading] at (12.05,16.5) {恢复的核心：D 补齐、G 对齐、R 从同一锚点重启};
\node[note,anchor=west,text width=23.0cm] at (.45,15.7) {以批次层故障为例：旧 G 到 200，D 到 150，共同前缀到 105；M 选择未来高度 $h_{\rm res}=260$。高度为压缩示意。};

\node[subheading,anchor=west] at (.45,14.7) {恢复前：旧生产分支与 D 已经不同；D 的已派生前缀必须保留};
\node[rowlabel] at (2.05,13.3) {旧 G / R};
\node[box,kept,fill=settled,text width=3.5cm] (oldbase) at (4.35,13.3) {\textbf{共同前缀}\\截至 105};
\node[box,dead,text width=5.2cm] (oldtail) at (10.85,13.3) {\textbf{旧 G：106--200 普通块}\\旧 R 包含同一旧后缀};
\draw[flow] (oldbase) -- (oldtail);
\node[note,anchor=west,text=rose,text width=7.0cm] at (14.3,13.3) {旧 session 失去兼容性时由 Manager Kill；\\旧 G 的后缀在 RESUME 时才被替换。};
\node[rowlabel,text=violet] at (2.05,11.35) {D};
\node[box,kept,fill=settled,text width=3.5cm] (dbase) at (4.35,11.35) {\textbf{共同前缀}\\截至 105};
\node[box,draw=violet,fill=native,text width=5.2cm] (dtail) at (10.85,11.35) {\textbf{106--150：deposit-only}\\窗口回退已确定的块};
\draw[flow,draw=violet] (dbase) -- (dtail);
\node[note,anchor=west,text=green,text width=7.0cm] at (14.3,11.35) {保留 D 的完整前缀和执行状态，\\恢复不改写这些已经确定的块。};

\node[card,draw=rose,fill=forceB,text width=21.5cm] at (12.05,8.95) {
  \textbf{M 提交 RESUME(260)，Authority 同一交易将 cursor override 到 260。}\\
  L1 最终化这笔交易后，派生先处理此前已接纳输入，再把空缺补到 260；GFR 的公共输入齐备后，采纳该次派生结果。
};

\node[subheading,anchor=west] at (.45,7.45) {恢复后：G 采用 D 截至 260 的前缀；新 R 从这个确定的块继续};
\node[rowlabel,text=violet] at (2.05,6.2) {D};
\node[rowlabel,text=green] at (2.05,4.2) {G};
\node[rowlabel,text=blue] at (2.05,2.2) {新 R};

\node[box,kept,fill=settled,text width=3.5cm] (d1) at (4.35,6.2) {\textbf{至 105}\\共同前缀};
\node[box,draw=violet,fill=native,text width=3.7cm] (d2) at (9.6,6.2) {\textbf{106--150}\\原有派生结果};
\node[box,draw=rose,fill=forceB,text width=3.7cm] (d3) at (14.85,6.2) {\textbf{151--260}\\deposit-only 补齐};
\draw[flow,draw=violet] (d1) -- (d2); \draw[flow,draw=rose] (d2) -- (d3);
\node[note,anchor=west,text width=5.4cm] at (17.7,6.2) {补到未来高度 260，\\留出状态准备时间。};

\node[box,kept,fill=settled,text width=3.5cm] (g1) at (4.35,4.2) {\textbf{至 105}\\保留共同前缀};
\node[box,draw=violet,fill=native,text width=3.7cm] (g2) at (9.6,4.2) {\textbf{106--150}\\采用 D 的同一批块};
\node[box,draw=rose,fill=forceB,text width=3.7cm] (g3) at (14.85,4.2) {\textbf{151--260}\\采用 D 的同一批块};
\draw[flow] (g1) -- (g2); \draw[flow] (g2) -- (g3);
\draw[ref,draw=rose] (d3.south) -- (g3.north);
\node[note,anchor=west,text=green,text width=5.4cm] at (17.7,4.2) {此时 G 的末块是 D[260]；\\旧 G 的 106--200 已替换。};

\node[box,kept,fill=settled,text width=14.2cm,minimum height=1.1cm] (inherit) at (9.6,2.2) {\textbf{新 session 的 anchor = D[260]}\\继承它的完整前缀与精确执行状态，创建新的 Program};
\draw[recover] (g3.south) -- (14.85,2.81);
\node[box,draw=blue,fill=regionA,text width=4.5cm,minimum height=1.1cm] (new) at (20.9,2.2) {\textbf{261、262\ldots}\\新 session 继续生产};
\draw[flow,draw=blue] (inherit) -- (new);

\node[note,anchor=west,text width=23.0cm] at (.45,.55) {正常生产仍要求区域 BFT 纳入全部应有的 deposit。恢复锚点固定为这次 RESUME 的 D[260]，不随节点后来看到的 D 末端移动。};
\node[note,anchor=west,text width=23.0cm] at (.45,-.3) {新 session 在后续 Espresso 块上完成分配并准备状态；此前已接纳数据若取不回，派生与 GFR 等待。};
\end{tikzpicture}
